\begin{longtblr}[
  entry=none,
  label={tab:d2_p1_fuzzy_approach_variant_notes},
  caption={Fase 1 (D2): variantes detectadas en D2\_Fuzzy\_approach\_type, propuesta de normalizacion y evidencia de soporte. Elaboracion propia.}
]{
  colspec={X[0.65,l,m] X[1.2,l,m] X[1.2,l,m] X[1.3,l,m] X[2.65,l,m]},
  width=\textwidth,
  row{1}={font=\bfseries}
}
Article\_ID & Variant\_detected & Proposed\_normalized & Evidence\_D2\_pages & Technical\_note \\
A135 & ADFIST & ADFIST & p.6-p.9 & Mantener separado de ANFIS: enfoque arbolado (Adaptive Dynamic FIS Tree) con optimizacion especifica reportada. \\
A144 & ADFIST & ADFIST & p.2-p.5 & Mantener separado de ANFIS: enfoque arbolado (Adaptive Dynamic FIS Tree) con optimizacion especifica reportada. \\
A168 & I-ANFIS & ANFIS & p.1; p.5-p.6; p.8 & Variante de la familia ANFIS; se normaliza para consolidar conteos por familia. \\
A157 & Other:Fuzzy\_ARTMAP & Hybrid\_ML\_Fuzzy & p.1-p.3 & Fuzzy ARTMAP integra componente neuro-fuzzy/ML; no corresponde a Mamdani/Sugeno puro. \\
A087 & Other:RANFIS & ANFIS & p.1; p.7-p.12 & Reportado como variante reforzada de ANFIS en el estudio. \\
A132 & Other:TSK & Sugeno/TSK & p.3; p.5-p.7 & TSK es subfamilia de Sugeno; se unifica para analisis comparativo. \\
A165 & Other:Tsukamoto & Tsukamoto & p.2-p.4 & Tsukamoto se mantiene como categoria propia (no equivalente a Sugeno/TSK). \\
A009 & Sugeno & Sugeno/TSK & p.2-p.3 & Se unifica con TSK para reducir dispersion semantica en la sintesis. \\
A123 & Sugeno & Sugeno/TSK & p.3-p.5 & Se unifica con TSK para reducir dispersion semantica en la sintesis. \\
A136 & Sugeno & Sugeno/TSK & p.12-p.16 & Se unifica con TSK para reducir dispersion semantica en la sintesis. \\
A165 & Type1\_FIS & Type1\_Generic & p.2-p.4 & Etiqueta generica tipo-1 sin subtipo explicito; se remapea para evitar ambiguedad. \\
A167 & Type1\_FIS & Type1\_Generic & p.1; p.3-p.4 & Etiqueta generica tipo-1 sin subtipo explicito; se remapea para evitar ambiguedad. \\
\end{longtblr}
